\section{DESENVOLVIMENTO TÉCNICO}

***Trata de apresentar a forma como foi desenvolvido o trabalho; O texto deve começar com a identificação do artefato, mencionando seu nome, o tipo (por exemplo: protótipo, documento, ferramenta, sistema, modelo, entre outros) e os autores envolvidos em sua criação.

Em seguida, deve-se apresentar o objetivo do artefato, explicando qual problema, necessidade ou desafio motivou sua elaboração e quais resultados se esperava alcançar com ele.

Após essa contextualização, deve-se detalhar o processo de desenvolvimento, incluindo as etapas seguidas, os métodos utilizados, as decisões tomadas, as ferramentas aplicadas, os testes realizados (quando houver) e as dificuldades encontradas ao longo do caminho. Também é importante justificar escolhas feitas durante o processo, apontando como elas contribuíram para a efetividade do artefato.